\section*{CONCLUSÕES}

Os ensaios experimentais sugerem que uma abordagem orientada a eventos pode auxiliar na integração de dados governamentais dispersos sobre queimadas. O patamar de $\sim$2,0 QPS sustentado com 2 vCPU frente à expressiva saturação com 1 vCPU denota que a antecipação de cálculos em memória favorece o tempo de resposta, embora demande dimensionamento cauteloso de hardware sob concorrência. A correlação observada na TI Parque do Araguaia indica o potencial desse modelo para complementar o suporte analítico à Defesa Civil, servindo como camada informativa adicional para o planejamento preventivo em campo.
